\chapter{Introducción y motivación del dominio de búsqueda} 
\label{chap:introduccion} 
Este capítulo presenta la motivación y los objetivos del proyecto.
Viajar permite conocer nuevos lugares, culturas, tradiciones y formas de vida.
Sin embargo, la preparación de un viaje puede requerir la consulta de numerosas 
páginas web para encontrar información sobre destinos, rutas, actividades o puntos 
de interés. A partir de esta necesidad se plantea el desarrollo de un buscador 
especializado en viajes por España, diseñado para reunir información turística 
procedente de distintas fuentes y facilitar su consulta.

\section{Motivación}
La elección del dominio está estrechamente relacionada con la composición y los 
intereses del grupo. Sus integrantes proceden de diferentes nacionalidades, lo que
 aporta una visión multicultural y distintas experiencias en torno a los viajes. 
 Al mismo tiempo, todos comparten un gran interés por descubrir lugares, conocer
  nuevos destinos y acceder a información que ayude a organizar futuras visitas.

Esta afinidad llevó al equipo a seleccionar un ámbito que resultase motivador y cercano,
 pero que también tuviese utilidad práctica. El turismo ofrece una gran variedad de 
 contenidos sobre ciudades, regiones, rutas, naturaleza, gastronomía y patrimonio. No 
 obstante, estos contenidos suelen estar repartidos entre diferentes portales y guías,
  por lo que un sistema de búsqueda específico puede ayudar a localizarlos de forma más rápida y ordenada.

De este modo, el proyecto combina la motivación personal de un equipo multicultural
interesado en viajar con el propósito académico de aplicar técnicas de recopilación,
tratamiento, indexación y búsqueda de información. El resultado pretende ser una herramienta
útil para explorar destinos españoles y obtener información relevante de una manera centralizada y accesible.

Algunos ejemplos de consulta que podrían realizarse con el buscador incluyen:
\begin{itemize}
    \item Lugares que visitar en Santiago de Compostela.
    \item Rutas de naturaleza en Andalucía.
    \item Gastronomía y enoturismo en La Rioja.
    \item Playas de las Islas Baleares.
\end{itemize}

\section{Objetivos} 
El objetivo general consiste en desarrollar un buscador especializado en información
turística sobre España. Los objetivos específicos son los siguientes: 
\begin{enumerate} 
\item Obtener documentos de diferentes fuentes mediante crawling. 
\item Limpiar y normalizar los datos recopilados. 
\item Indexar los documentos mediante Elasticsearch. 
\item Diseñar consultas adaptadas al dominio turístico. 
\item Desarrollar una interfaz web para realizar búsquedas. 
\end{enumerate} 
\section{Estructura de la memoria} 
La memoria se organiza en cinco capítulos. Después de esta introducción, se presentan
el dominio y las fuentes de datos. A continuación, se describe el desarrollo del sistema,
incluyendo el crawling, la preparación de los datos y la indexación. Finalmente, se explican
las decisiones de búsqueda y las funcionalidades de la aplicación.